% -*- mode:LaTex; mode:visual-line; mode:flyspell; fill-column:75-*-

\chapter{Background} \label{chapBackground}

% TODO: Brief chapter intro — what this chapter covers and how it
% motivates the technical choices made in later chapters.

\section{Agricultural Robotics} \label{secAgRobotics}

\subsection{Agricultural Robots Today} \label{secAgRoboticsToday}

% TODO: Survey robotic platforms and systems deployed in agricultural
% settings today.
% - Harvesting robots, pruning robots, scouting/inspection robots.
% - Challenges: unstructured environments, occlusion, lighting variation,
%   terrain.
% - Relevant platforms: ground robots (Amiga, Thorvald, etc.),
%   manipulator arms for targeted inspection.
% - Cite key related systems; compare to your Amiga + xArm6 setup.

\section{Fire Blight Disease} \label{secFireBlight}

% TODO: Brief intro to fire blight as the disease motivating this thesis.

\subsection{Why Fire Blight Matters} \label{secWhyFireBlight}

% TODO: Motivate the economic and agronomic importance of fire blight.
% - Crops affected (apple, pear, other Rosaceae).
% - Economic impact: cite orchard loss statistics.
% - Speed of spread and risk of losing entire orchards/blocks.

\subsection{Disease Mechanism} \label{secFireBlightMechanism}

% TODO: Biology and epidemiology of Erwinia amylovora.
% - Infection mechanism: enters via blossoms and wounds, spreads through
%   vascular tissue.
% - Visible symptoms and progression: shepherd's crook wilt,
%   water-soaked lesions, orange/brown discoloration of shoots and fruit.
% - Detection challenges: early infection may not be visually obvious;
%   NIR reflectance changes may precede visible symptoms.

\subsection{Current Solutions and Problems} \label{secFireBlightCurrentSolutions}

% TODO: Describe how orchardists currently detect and respond to fire
% blight, and the shortcomings of these approaches.
% - Manual scouting: labor-intensive, infrequent, error-prone.
% - Management practices: pruning, copper/antibiotic sprays.
% - Why current practice motivates automated, systematic detection.

\section{Current Efforts on Automated Fire Blight and Orchard Disease Diagnosis} \label{secCurrentEfforts}

% TODO: Brief intro framing the gap this thesis addresses.

\subsection{Datasets} \label{secCurrentEffortsDatasets}

% TODO: Review existing public/research datasets for fire blight and
% orchard disease detection.
% - Modalities used (RGB, hyperspectral/NIR).
% - Scale and limitations (size, diversity, label quality).
% - How your dataset (Chapter~\ref{chapDataset}) compares/improves on these.

\subsection{Orchard Robots} \label{secCurrentEffortsRobots}

% TODO: Review existing robotic/automated systems for orchard disease
% inspection and diagnosis.
% - Sensing approaches (RGB, hyperspectral, LiDAR), autonomy level,
%   reported performance.
% - Gaps: lack of 3D semantic mapping, lack of active/NBV-driven
%   inspection strategies.

\section{Summary} \label{secBackgroundSummary}

% TODO: Summarize what gaps in the literature your thesis addresses.
% Tie each gap to a contribution stated in Chapter~\ref{chapIntroduction}.
